% Compiler avec pdfLaTeX :  pdflatex lebenslauf_muster.tex  (2 fois)

\documentclass[11pt,a4paper]{article}

\usepackage[T1]{fontenc}
\usepackage[utf8]{inputenc}
\usepackage[ngerman]{babel}
\usepackage[scaled]{helvet}
\renewcommand{\familydefault}{\sfdefault}
\usepackage{microtype}
\usepackage[a4paper,left=1.8cm,right=1.8cm,top=1.2cm,bottom=1.5cm]{geometry}
\usepackage{xcolor}
\usepackage{tikz}
\usepackage{enumitem}
\usepackage[hidelinks]{hyperref}

\definecolor{shape}{HTML}{4F6FC0}
\definecolor{accent}{HTML}{5B7DB1}
\definecolor{textgray}{HTML}{3A3A3A}
\definecolor{lightgray}{HTML}{8A8A8A}

\pagestyle{empty}
\setlength{\parindent}{0pt}
\color{textgray}

% ---------- Commandes ----------
\newcommand{\sectionhead}[1]{%
  \vspace{1.1em}
  {\large\bfseries\color{accent}\MakeUppercase{#1}}\\[-0.35em]
  {\color{accent!50}\rule{\linewidth}{0.5pt}}\vspace{0.5em}}

% \entry{date}{titre}{organisation}{puces}
\newcommand{\entry}[4]{%
  \noindent
  \begin{minipage}[t]{2.9cm}\raggedright\small\color{lightgray}#1\end{minipage}%
  \hspace{0.3cm}%
  \begin{minipage}[t]{\dimexpr\linewidth-3.2cm\relax}
    {\normalsize\bfseries\color{black!90}#2}\\[-0.1em]
    {\small\color{accent}#3}%
    \if\relax\detokenize{#4}\relax
    \else
      \begin{itemize}[leftmargin=1.1em,itemsep=0.1em,topsep=0.3em,parsep=0pt,label={\tiny$\bullet$}]
        \small #4
      \end{itemize}
    \fi
  \end{minipage}\par\vspace{0.55em}}

\begin{document}

% ---------- Kopfbereich ----------
\begin{tikzpicture}[remember picture,overlay]
  % Dekoratives blaues Element oben
  \fill[shape!85,rounded corners=8pt]
    ([xshift=7.0cm]current page.north west) --
    ([xshift=13.0cm]current page.north west) --
    ([xshift=11.0cm,yshift=-1.4cm]current page.north west) --
    ([xshift=8.0cm,yshift=-0.9cm]current page.north west) -- cycle;
  % Name
  \node[anchor=north west,inner sep=0pt] at ([xshift=1.8cm,yshift=-2.0cm]current page.north west)
    {\fontsize{24}{28}\selectfont\bfseries\color{black!80}Emmanuel Dakdem Nguena};
  % Kontakt
  \node[anchor=north east,inner sep=0pt,align=right,font=\footnotesize,text=black!65]
    at ([xshift=-1.8cm,yshift=-1.0cm]current page.north east)
    {Ederstraße 6, 35390 Gießen\\
     017643614045\\
     emmanuel.dakdem@yahoo.com\\
     Führerschein B197};
\end{tikzpicture}

\vspace*{2.7cm}

% Profil
{\small\linespread{1.08}\selectfont
Maschinenbau- und Mechatronik-Student mit starkem Fokus auf Konstruktion, CAD-basierter Produktentwicklung
und technischer Dokumentation. Durch die Bachelorarbeit und praktische Projekte mit Siemens NX, Ansys und
SolidWorks habe ich fundierte Erfahrungen in der Entwicklung technischer Baugruppen, der Erstellung von
Konstruktionsunterlagen und der fachübergreifenden Zusammenarbeit in industriellen Prozessen. Meine
Stärken liegen in strukturierter Arbeit, Prozessverständnis und der Umsetzung technischer Lösungen mit
hohem Qualitätsanspruch.\par}

\sectionhead{Berufserfahrung}

\entry{Apr. 2024 --\\ Feb. 2025}
{Werkstudent neue Technologien und Konnektivität -- Aßlar}
{PFEIFFER VACUUM}
{\item Unterst"{u}tzung bei der SAP S/4HANA-Einf"{u}hrung und Abstimmung mit bestehenden Prozessen
 \item Strukturierte Dokumentation von Anforderungen und fach"{u}bergreifende Projektarbeit
 \item Koordination mit internen Teams und potenziellen externen Lieferanten
 \item Konzeptentwicklung zur kamerabasierten Qualit"{a}tssicherung bei einer Laserentgratanlage
 \item Datenrecherchen, Tool-Scouting und Unterst"{u}tzung bei kleineren technischen Projekten}

\entry{Apr. 2023 --\\ Apr. 2024}
{Werkstudent Qualit"{a}tskontrolle -- Heuchelheim}
{Schunk Carbon Technology GmbH}
{\item Durchf"{u}hrung von Messprotokollen und Dokumentation
 \item Analyse und Dokumentation von Produktionsdaten mit SAP
 \item Qualit"{a}tsorientierte Prozessunterst"{u}tzung und Fehleranalyse in der Produktion}

\sectionhead{Ausbildung}

\entry{Okt. 2025 --\\ heute}
{Master of Science (Schwerpunkt: Mechatronik \& Robotik) -- Friedberg}
{Technische Hochschule Mittelhessen}
{\item Vertiefung in interdisziplin"{a}ren Bereichen der Mechatronik und Robotik
 \item Fokus auf Mechanik, Elektronik, Informatik, Systemsimulation und KI-Methoden}

\entry{Okt. 2021 --\\ Sep. 2025}
{Bachelor of Engineering (Schwerpunkt: Allgemeiner Maschinenbau) -- Gießen}
{Technische Hochschule Mittelhessen}
{\item Konstruktion und Inbetriebnahme technischer Maschinen und Anlagen
 \item CAD- und Simulationssoftware (Siemens NX, Ansys, SolidWorks), Automatisierungs- und Messtechnik
 \item Fundierte Kenntnisse in Fertigungstechnologien, Werkstoffkunde und technischer Dokumentation}

\entry{Sep. 2016 --\\ Aug. 2018}
{Allgemeine Hochschulreife (Abitur) -- Yaoundé, Kamerun}
{Gymnasium CSAO}
{\item Naturwissenschaftlicher Schwerpunkt: Mathematik und Physik
 \item Fundierte Kenntnisse in analytischem Denken und technischer Probleml"{a}sung}

\sectionhead{Praktika \& Projekte}

\entry{Apr. 2026 --\\ heute}
{Projekt: Klassische und visuelle Odometrieverfahren f"{u}r ein automatisiertes Modellfahrzeug}
{Technische Hochschule Mittelhessen, Friedberg}
{\item \textbf{Konzeption \& Implementierung:} Modulares Odometriesystem in ROS~2 f"{u}r das Modellfahrzeug MxCarkit
 \item \textbf{Systemintegration:} Einbindung in eine bestehende ROS-2-Architektur mit standardisierten Schnittstellen
 \item \textbf{Validierung \& Analyse:} Experimentelle Fahrversuche zur Bewertung von Genauigkeit und Robustheit}

\entry{März 2025 --\\ Mai 2025}
{Praktikant: Analysetool zur Identifikation von Optimierungspotentialen in der Produktion}
{PFEIFFER VACUUM, Aßlar}
{\item Visualisierung von Produktionskennzahlen mit Power~BI und DAX
 \item Aufbereitung von Daten aus dem Dataflow, Erstellung von Dashboards und Berichten
 \item Analyse von Fertigungsprozessen und Identifikation von Engp"{a}ssen}

\sectionhead{Kompetenzen}

\noindent\small\textbf{Sprachen:} Französisch (verhandlungssicher), Deutsch (verhandlungssicher), Englisch (verhandlungssicher)\par
\vspace{0.4em}
\noindent\small\textbf{Software:} Siemens NX (CAD, FEM), SolidWorks, Ansys, SAP S/4HANA, Python, MS Office, Codesys, Siemens TIA Portal, ROS~2, C++\par

\end{document}
